%==============================================================
%  lecons/combinatoire/bell.tex — Nombres de Bell
%==============================================================

\lecontitle{Nombres de Bell}{Combinatoire — Partitions d'ensembles}

%=============================================================
%  § 1 — DÉFINITIONS
%=============================================================
\secnum{navyblue}{1}{Définitions et premières valeurs}

\begin{tcolorbox}[thmbox]
  \textbf{Définition.}\enspace\index{nombre de Bell}\index{partition d'un ensemble}
  Une \emph{partition}\index{partition} de l'ensemble $[n] = \{1,2,\ldots,n\}$
  est une famille de parties non vides, deux à deux disjointes, dont la réunion
  est $[n]$.

  \medskip
  Le \emph{$n$-ième nombre de Bell}, noté $B_n$, est le nombre de partitions
  de $[n]$. Par convention, $B_0 = 1$ (l'ensemble vide admet exactement une
  partition : la partition vide).
\end{tcolorbox}

\rubriq{Premières valeurs}

\begin{mdframed}[style=exbar]
  \textbf{$B_0 = 1$.}\enspace La seule partition de $\emptyset$ est $\{\}$.

  \smallskip\noindent
  \textbf{$B_1 = 1$.}\enspace La seule partition de $\{1\}$ est $\bigl\{\{1\}\bigr\}$.

  \smallskip\noindent
  \textbf{$B_2 = 2$.}\enspace Les partitions de $\{1,2\}$ sont
  $\bigl\{\{1,2\}\bigr\}$ et $\bigl\{\{1\},\{2\}\bigr\}$.

  \smallskip\noindent
  \textbf{$B_3 = 5$.}\enspace Les partitions de $\{1,2,3\}$ sont :
  $\{123\}$,
  $\{1\}\{23\}$, $\{2\}\{13\}$, $\{3\}\{12\}$,
  $\{1\}\{2\}\{3\}$.

  \medskip
  \begin{center}
    \begin{tabular}{c|ccccccc}
      $n$ & 0 & 1 & 2 & 3 & 4 & 5 & 6 \\\hline
      $B_n$ & 1 & 1 & 2 & 5 & 15 & 52 & 203
    \end{tabular}
  \end{center}
\end{mdframed}

\decsep{}

%=============================================================
%  § 2 — RÉCURRENCE ET TRIANGLE DE BELL
%=============================================================
\secnum{navyblue}{2}{Récurrence et triangle de Bell}

\begin{tcolorbox}[thmbox]
  \textbf{Théorème (récurrence de Bell).}\enspace\index{récurrence de Bell}
  Pour tout entier $n \geq 0$ :
  \[
    B_{n+1} = \sum_{k=0}^{n} \binom{n}{k} B_k.
  \]
\end{tcolorbox}

\rubriq{Preuve}

\begin{mdframed}[style=proofbar]
  On conditionne par le bloc contenant l'élément $n+1$.
  Ce bloc contient $n+1$ et $k$ autres éléments parmi $\{1,\ldots,n\}$
  ($0 \leq k \leq n$) : il y a $\binom{n}{k}$ façons de choisir ces $k$
  éléments. Les $n-k$ éléments restants se partitionnent en $B_{n-k}$ façons,
  indépendamment. Ainsi :
  \[
    B_{n+1} = \sum_{k=0}^{n} \binom{n}{k} B_{n-k}
             = \sum_{j=0}^{n} \binom{n}{n-j} B_j
             = \sum_{j=0}^{n} \binom{n}{j} B_j,
  \]
  où l'avant-dernière égalité s'obtient en posant $j=n-k$, et la dernière
  en utilisant $\binom{n}{n-j}=\binom{n}{j}$.
  $\blacksquare$
\end{mdframed}

\rubriq{Triangle de Bell}

\begin{mdframed}[style=exbar]
  \index{triangle de Bell}
  Le \emph{triangle de Bell} est un tableau triangulaire $(a_{n,k})_{0\leq k\leq n}$
  construit comme suit :
  \begin{itemize}
    \item $a_{n,0} = B_n$ (premier terme de chaque ligne).
    \item $a_{n,k} = a_{n,k-1} + a_{n-1,k-1}$ (chaque terme est la somme du
          terme précédent sur la même ligne et du terme au-dessus à gauche).
  \end{itemize}
  Le dernier terme de chaque ligne est $B_{n+1}$.

  \medskip
  \begin{center}
    \begin{tabular}{ccccc}
      1 & & & & \\
      1 & 2 & & & \\
      2 & 3 & 5 & & \\
      5 & 7 & 10 & 15 & \\
      15 & 20 & 27 & 37 & 52 \\
    \end{tabular}
  \end{center}
\end{mdframed}

\decsep{}

%=============================================================
%  § 3 — NOMBRES DE STIRLING ET LIEN AVEC $B_n$
%=============================================================
\secnum{navyblue}{3}{Nombres de Stirling de deuxième espèce}

\begin{tcolorbox}[thmbox]
  \textbf{Définition.}\enspace\index{nombre de Stirling!de deuxième espèce}
  Le \emph{nombre de Stirling de deuxième espèce} $S(n,k)$ est le nombre de
  partitions de $[n]$ en exactement $k$ blocs non vides.

  \medskip
  \textbf{Récurrence.}\enspace\index{récurrence de Stirling}
  \[
    S(n,k) = k\,S(n-1,k) + S(n-1,k-1),
    \qquad S(n,0)=[n=0],\quad S(n,n)=1.
  \]

  \medskip
  \textbf{Lien avec les nombres de Bell.}\enspace
  \[
    B_n = \sum_{k=0}^{n} S(n,k).
  \]
\end{tcolorbox}

\rubriq{Preuve de la récurrence de Stirling}

\begin{mdframed}[style=proofbar]
  Pour partitionner $[n]$ en $k$ blocs, on considère l'élément $n$.
  \begin{itemize}
    \item Soit $\{n\}$ forme un bloc seul : les $n-1$ éléments restants
          se partitionnent en $k-1$ blocs — $S(n-1,k-1)$ façons.
    \item Soit $n$ rejoint l'un des $k$ blocs d'une partition de $[n-1]$ :
          il y a $k\cdot S(n-1,k)$ façons.
  \end{itemize}
  D'où $S(n,k) = k\,S(n-1,k) + S(n-1,k-1)$. $\blacksquare$
\end{mdframed}

\rubriq{Formule explicite}

\begin{tcolorbox}[thmbox]
  \textbf{Formule d'inclusion-exclusion.}\enspace
  \[
    S(n,k) = \frac{1}{k!}\sum_{j=0}^{k}{(-1)}^{k-j}\binom{k}{j}j^n.
  \]
\end{tcolorbox}

\decsep{}

%=============================================================
%  § 4 — FORMULE DE DOBI\'NSKI ET SÉRIE GÉNÉRATRICE
%=============================================================
\secnum{navyblue}{4}{Formule de Dobi\'nski et série génératrice}

\begin{tcolorbox}[thmbox]
  \textbf{Formule de Dobi\'nski.}\enspace%
  \index{Dobi\'nski (formule de)}
  \[
    B_n = \frac{1}{e}\sum_{k=0}^{+\infty}\frac{k^n}{k!}.
  \]

  \medskip
  \textbf{Série génératrice exponentielle.}\enspace%
  \index{série génératrice exponentielle!nombres de Bell}
  \[
    \sum_{n=0}^{+\infty} B_n\,\frac{x^n}{n!} = e^{e^x - 1}.
  \]
\end{tcolorbox}

\rubriq{Preuve de la formule de Dobi\'nski}

\begin{mdframed}[style=proofbar]
  On part de la formule d'inclusion-exclusion pour $S(n,k)$ et on somme sur
  $k$. En écrivant $\binom{k}{j}=\dfrac{k!}{j!\,(k-j)!}$, le facteur $1/k!$ se
  simplifie :
  \[
    B_n = \sum_{k=0}^{n} S(n,k)
        = \sum_{k=0}^{n}\frac{1}{k!}\sum_{j=0}^{k}{(-1)}^{k-j}\binom{k}{j}j^n
        = \sum_{k=0}^{n}\sum_{j=0}^{k}\frac{{(-1)}^{k-j}}{j!\,(k-j)!}\,j^n.
  \]
  On peut prolonger la sommation sur $k$ jusqu'à $+\infty$ sans rien changer :
  pour $k>n$, il n'existe aucune partition de $[n]$ en $k$ blocs, donc
  $S(n,k)=0$ et les termes ajoutés sont nuls. Dans la double somme ainsi
  obtenue, seuls les indices $0\leq j\leq k$ contribuent.
  Intervertissons alors les deux sommations (la série double est absolument
  convergente, $j^n/(j!\,(k-j)!)\geq 0$), en sommant d'abord sur $k\geq j$ :
  \[
    B_n = \sum_{j=0}^{+\infty}\frac{j^n}{j!}
          \sum_{k\geq j}\frac{{(-1)}^{k-j}}{(k-j)!}.
  \]
  Le facteur intérieur ne dépend plus que de $k-j$ : en posant $\ell=k-j$,
  \[
    \sum_{k\geq j}\frac{{(-1)}^{k-j}}{(k-j)!}
      = \sum_{\ell=0}^{+\infty}\frac{{(-1)}^{\ell}}{\ell!}
      = e^{-1}.
  \]
  En reportant, il vient
  \[
    B_n = e^{-1}\sum_{j=0}^{+\infty}\frac{j^n}{j!}.
  \]
  $\blacksquare$
\end{mdframed}

\rubriq{Preuve de la série génératrice}

\begin{mdframed}[style=proofbar]
  On dispose, via la formule de Dobi\'nski, de :
  \[
    B_n = e^{-1}\sum_{k\geq 0}\frac{k^n}{k!}.
  \]
  Ainsi :
  \[
    \sum_{n\geq 0}B_n\frac{x^n}{n!}
    = e^{-1}\sum_{k\geq 0}\frac{1}{k!}\sum_{n\geq 0}\frac{{(kx)}^n}{n!}
    = e^{-1}\sum_{k\geq 0}\frac{e^{kx}}{k!}
    = e^{-1}\,e^{e^x}
    = e^{e^x - 1}.
  \]
  $\blacksquare$
\end{mdframed}

\decsep{}

%=============================================================
%  § 5 — APPLICATIONS, PIÈGES, EXERCICES
%=============================================================
\secnum{exgreen}{5}{Applications, pièges et exercices}

\rubriq{Applications}

\begin{mdframed}[style=exbar]
  \textbf{1. Classement de structures.}\enspace%
  \index{relation d'équivalence}
  Le nombre de relations d'équivalence sur $[n]$ est $B_n$
  (une relation d'équivalence = une partition en classes).

  \smallskip\noindent
  \textbf{2. Moments d'une loi de Poisson.}\enspace%
  \index{loi de Poisson}
  Si $X \sim \mathrm{Poisson}(1)$, alors $\mathbb{E}[X^n] = B_n$.
  Plus généralement, si $X \sim \mathrm{Poisson}(\lambda)$,
  \[
    \mathbb{E}[X^n] = \sum_{k=0}^{n} S(n,k)\,\lambda^k.
  \]

  \smallskip\noindent
  \textbf{3. Dérivation de fonctions composées.}\enspace%
  \index{formule de Fa\`a di Bruno}
  La formule de Fa\`a di Bruno exprime ${(f \circ g)}^{(n)}$ en termes
  des dérivées de $f$ et $g$ à l'aide des \emph{polynômes de Bell partiels}
  $B_{n,k}\bigl(g',g'',\ldots,g^{(n-k+1)}\bigr)$ :
  \[
    {(f\circ g)}^{(n)}
      = \sum_{k=0}^{n} f^{(k)}\!\circ g \;\cdot\;
        B_{n,k}\bigl(g',\ldots,g^{(n-k+1)}\bigr).
  \]
  (Les nombres de Stirling de deuxième espèce s'en déduisent comme cas
  particulier : $S(n,k)=B_{n,k}(1,1,\ldots,1)$.)
\end{mdframed}

\rubriq{Pièges}

\begin{mdframed}[style=cexbar]
  \textbf{$B_n$ croît plus vite que toute exponentielle fixée.}\enspace%
  On a $B_n \geq S(n,1) = 1$ et $B_n \geq S(n,n) = 1$, mais la croissance
  réelle est bien plus rapide :
  \[
    \ln B_n = n\bigl(\ln n - \ln\ln n - 1\bigr) + o(n)
  \]
  (asymptotique de Moser--Wyman / de Bruijn). On confondra $B_n$ avec $n!$ ou
  avec $2^n$ par erreur.

  \smallskip\noindent
  \textbf{Confondre $S(n,k)$ avec $\begin{bmatrix}n\\k\end{bmatrix}$.}\enspace%
  Les \emph{nombres de Stirling de première espèce}
  $\begin{bmatrix}n\\k\end{bmatrix}$ comptent les permutations de $[n]$ à
  $k$ cycles, et non les partitions. Ils vérifient une récurrence similaire
  mais avec un signe différent.

  \smallskip\noindent
  \textbf{Convention $B_0 = 1$.}\enspace%
  L'ensemble vide admet exactement une partition (la partition vide). Ne pas
  confondre avec $B_1 = 1$ : les deux valent $1$, mais pour des raisons
  différentes.
\end{mdframed}

\vspace{6pt}
\exolabel

\begin{mdframed}[style=exobar]
  \begin{enumerate}
    \item En utilisant la récurrence $B_{n+1} = \sum_{k=0}^{n}\binom{n}{k}B_k$,
          calculer $B_6 = 203$ à partir des valeurs précédentes.

    \item Montrer que $S(n,2) = 2^{n-1}-1$ pour $n\geq 1$.
          (\textit{Indication :} une partition en 2 blocs correspond au choix du
          bloc contenant $1$, qui ne peut pas être $[n]$ tout entier.)

    \item Vérifier sur $n=3$ que $\mathbb{E}[X^3] = B_3 = 5$ pour
          $X\sim\mathrm{Poisson}(1)$ en calculant directement
          $\sum_{k\geq 0} k^3\,e^{-1}/k!$.

    \item Montrer que $B_n$ est impair si et seulement si $n\equiv 0$ ou
          $1\pmod{3}$ (\textit{i.e.} $B_n$ est pair si et seulement si
          $n\equiv 2\pmod{3}$). La suite $(B_n \bmod 2)$ est en effet
          périodique de période $3$ : $1,1,0,1,1,0,\ldots$

    \item (\textit{Défi}) En utilisant la série génératrice $e^{e^x-1}$,
          retrouver la formule de Dobi\'nski en développant $e^{e^x}$ comme
          double somme.
  \end{enumerate}
\end{mdframed}

\leconfooter{%
  É.\ Dobi\'nski (1877) · E.\ T.\ Bell, \textit{Am.\ Math.\ Monthly} (1934)%
}
